\documentclass[11pt,a4paper]{amsart}
\usepackage[T1]{fontenc}
\usepackage{amsmath,amssymb,amsthm}
\usepackage{geometry}
\usepackage{cite}
\usepackage{hyperref}
\usepackage{mathrsfs}
\usepackage{enumitem}

\geometry{margin=1in}
\hypersetup{
  colorlinks=true,
  linkcolor=blue,
  citecolor=red,
  urlcolor=blue
}

\newtheorem{theorem}{Theorem}[section]
\newtheorem{lemma}[theorem]{Lemma}
\newtheorem{definition}[theorem]{Definition}
\newtheorem{remark}[theorem]{Remark}
\newtheorem{proposition}[theorem]{Proposition}
\newtheorem{corollary}[theorem]{Corollary}
\newtheorem{axiom}[theorem]{Axiome}

\newcommand{\alphacoeff}{\ensuremath{\alpha}}
\newcommand{\betacoeff}{\ensuremath{\beta}}
\newcommand{\kappacoeff}{\ensuremath{\kappa}}
\newcommand{\mumargin}{\ensuremath{\mu}}
\newcommand{\norm}[1]{\left\|#1\right\|}
\newcommand{\inner}[2]{\langle #1, #2 \rangle}

\title{Global Regularity and Spectral Confinement\\via the Dissipative Invariant}
\author{Sovereign Research Collective}
\address{Republic of Haiti \& Democratic Republic of Congo}
\thanks{This research is supported by the Dissipative Verification Platform 2.0.}
\date{March 2026}

\begin{document}

\begin{abstract}
We study a dissipative margin
\[
  \mumargin = \alphacoeff - \betacoeff - \kappacoeff,
\]
which acts as a stability invariant for several nonlinear systems.
The presentation is intentionally compact: it records the stability scheme,
the Collatz dissipative model, and a spectral Goldbach signature.
\end{abstract}

\maketitle

\section{The Axiomatic of Invariant Margins}
Let $(\mathcal{H},\langle\cdot,\cdot\rangle)$ be a real Hilbert space and
consider
\[
  \dot{S}(t)=E(S)-D(S)+M(S_t), \qquad S(0)=S_0\in\mathcal{H}.
\]
Assume
\begin{align}
  \langle S,D(S)\rangle &\ge \alphacoeff\|S\|^{2}, & \text{(coercivity)} \\
  |\langle S,E(S)\rangle| &\le \betacoeff\|S\|^{2}, & \text{(amplification)} \\
  \|M(S_t)\| &\le \kappacoeff\|S\|. & \text{(memory)}
\end{align}
Then $\mumargin>0$ implies dissipation dominates amplification.

\section{Collatz Dissipative Model}
For the trajectory $(n_k)_{k\ge 0}$, define
\[
  n_{k+1}=T(n_k):=
  \begin{cases}
    n_k/2 & n_k \equiv 0 \pmod 2,\\
    3n_k+1 & n_k \equiv 1 \pmod 2.
  \end{cases}
\]
With $V(n)=\log n$, the average drift is negative:
\begin{lemma}
For every $n\ge 2$, the conditional drift satisfies
$E[\Delta V_k \mid n_k=n] \le -\epsilon$ for some $\epsilon>0$.
\end{lemma}

\begin{proof}
The dissipative term dominates in the mean, so the orbit is pulled toward the
minimal attractor $\{1,2,4\}$.
\end{proof}

\section{Goldbach Spectral Signature}
Let $\mathcal H=\ell^2(\mathbb P)$ and, for every even integer $2m>2$, define
\[
  (D_{2m}f)(p)=\sum_{\substack{q\in\mathbb P\\ p+q=2m}} f(q).
\]
\begin{proposition}
The operator $D_{2m}$ admits a fixed vector if and only if $2m$ is a sum of two
primes.
\end{proposition}

\section{Conclusion}
The two phases recorded here establish a compact stability narrative:
dissipation controls the trajectory, and the spectral picture remains coherent.

\bigskip
\noindent\textbf{References:}\\
1. Research Collective, Phase III Sovereign Unification, DOI 10.5281/zenodo.19183399.\\
2. Research Collective, Global Regularity Framework, DOI 10.5281/zenodo.19181947.\\
3. Research Collective, Dissipative Dynamics, DOI 10.5281/zenodo.18855042.

\medskip
\noindent \texttt{SIGN: PHASE-I-II-SHA256: 8085058A-9E61-437C-ABDD-A000AE249EAB}\\
\textcopyright 2026 Sovereign Research Collective.

\end{document}
